\SourceLeaf{407}%
\[
\text{V.}\quad .\ .\ .\ .\ .\ .\ .\qquad b=90^0,\quad A\gtrless90^0
\]
\par
1) \(c=90^0\). Dann wird \(A\) der Pol von \(c\), also \(b=90^0\), und \(a=A\). Die Gleichung
\[
\cos a=\cos b.\cos c+\sin b.\sin c.\cos A
\]
ist von selbst evident.

2) \(c<90^0\). Hier wird \(c\) über \(B\) hinaus bis zur Länge von \(90^0\) fortgesetzt. Aus der Betrachtung des Dreiecks folgt dann
\[
\cos a=\cos(90^0-c).\cos A+\sin(90^0-c).\sin A.\cos90^0
\]
oder
\[
\cos a=\sin c.\cos A
\]
daher die Formel auch in diesem Falle richtig ist.

3) \(c>90^0\). Hier wird von \(c\) der Bogen \(90^0\) in dem Punkte \(R\) abgeschnitten, dann folgt aus Betrachtung des Dreiecks \(BRC\)
\[
\cos a=\cos(c-90^0).\cos A+\sin(c-90^0).\sin A.\cos90^0
\]
oder
\[
\cos a=\sin c.\cos A
\]
wie im vorigen Fall.

Die Fälle, wo \(b<90^0\) oder \(b>90^0\), und zugleich \(c=90^0\), sind mit den beiden vorigen wesentlich einerlei.

Zählt man also alle möglichen Fälle auf, so folgt der Beweis, wenn
\begin{center}
\begin{tabular}{ll}
Beide Seiten kleiner als \(90^0\)&aus I\\
Beide Seiten grösser als \(90^0\)&aus II\\
Eine grösser, die andere kleiner als \(90^0\)&aus III\\
Beide \(=90^0\)&aus IV und V.\,1.\\
Eine \(=90^0\), die andere kleiner&aus IV und V.\,2.\\
Eine \(=90^0\), die andere grösser&aus IV und V.\,3.
\end{tabular}
\end{center}
\hspace*{\parindent}Wir können also allgemein annehmen:
\[
\text{(A)}\quad .\ .\ .\ .\ .\quad\cos a=\cos b.\cos c+\sin b.\sin c.\cos A
\]
und ebenso
